\chapter{Evaluation}

This chapter presents the experimental results, hypothesis test outcomes, and interpretation of findings in the context of the research question and objectives.

\section{Experimental Execution Summary}

A total of \textbf{350 experimental runs} were executed in the local simulator across five phases:

\begin{table}[h]
\centering
\begin{tabular}{lrr}
\hline
\textbf{Phase} & \textbf{Runs} & \textbf{Purpose} \\
\hline
Pilot & 18 & Repeatability validation, power analysis \\
Baseline Replication & 115 & Steady-state comparison with \citet{Kyrychenko2025} \\
Arms Comparison & 10 & Sync vs. queue under no fault and one fault \\
Fault Campaigns A-E & 187 & Systematic VT × MRC × fault grid for H1-H3 \\
Burst Replication & 20 & Key cells under higher load \\
\hline
\textbf{Total} & \textbf{350} & \\
\hline
\end{tabular}
\caption{Experimental phases and run counts}
\end{table}

All runs completed successfully; no crashes or invalid manifests. Manifests are committed to \texttt{results/<phase>/manifests/*.json}. Figures are in \texttt{results/figures/*.png}.

\section{Pilot Phase: Repeatability and Power}

\subsection{Repeatability}

The pilot tested 3 configurations (VT 30/300/900, batch 10, MRC 3, no fault) with 3 repeats each. Coefficient of variation (CV) for key metrics:
\begin{itemize}
    \item Throughput: CV < 0.01 (highly repeatable)
    \item Latency: CV < 0.05 (acceptable)
    \item Loss: 0.0 in all runs (no variance)
    \item Duplicates: CV ~0.2 (expected for stochastic at-least-once)
\end{itemize}

\textbf{Conclusion:} Simulator produces consistent results; seeds control randomness effectively.

\subsection{Power Analysis}

Pilot power planning (`power_from_pilot.json`) recommended 3 confirmatory repeats (plan\_repeats=5). Committed confirmatory cells in \texttt{stats\_H1\_H2\_H3.json} use $n=10$ per group.

\section{Baseline Replication: Steady-State Validation}

\subsection{Comparison with Kyrychenko et al. (2025)}

The baseline phase replicated \citet{Kyrychenko2025} steady-state conditions (no faults). Key findings:

\begin{itemize}
    \item \textbf{Batch size dominates throughput:} Median throughput ~161 msg/s at batch 1, ~375 msg/s at batch 50 (5 pollers). Matches Kyrychenko's qualitative finding ("batch size is key factor").
    
    \item \textbf{Visibility timeout does not affect steady-state throughput:} Throughput variance across VT 30-900 < 5\%. Confirms Kyrychenko's result.
    
    \item \textbf{Duplicate floor:} Median duplicate rate ~0.002 (0.2\%) from at-least-once delivery, consistent across configs. No simulation artifact.
    
    \item \textbf{Delivery delay shifts latency:} Delay 300 adds ~300s to median latency; throughput unaffected. Expected.
\end{itemize}

\textbf{Conclusion:} Simulator replicates steady-state \emph{relative effects} from Kyrychenko et al. Absolute throughput differs (simulator is faster, no network overhead), but configuration trade-offs match. Simulator is valid for testing fault scenarios.

\section{Arms Comparison: Control Validation}

The sync vs. queue arms were compared under (1) no fault and (2) consumer\_kill fault (10s).

\textbf{No Fault:}
\begin{itemize}
    \item Sync arm: 0.0\% loss, median latency ~50ms, throughput limited by sequential processing.
    \item Queue arm: 0.0\% loss, median latency ~120ms (includes queue delay), throughput ~350 msg/s (batch 50).
\end{itemize}

\textbf{Consumer Kill (10s):}
\begin{itemize}
    \item Sync arm: 0.3\% loss (client retry budget exhausted for messages sent during fault window), recovery immediate after fault stop.
    \item Queue arm: 0.0\% loss (redelivery from SQS), recovery time ~35s (VT 30).
\end{itemize}

\textbf{Interpretation:} Sync arm demonstrates that faults \emph{do} cause loss without queueing; queue arm's zero loss validates DLQ + redelivery mechanism. Control arm is valid baseline.

\section{Hypothesis Test Results}

All hypothesis tests used α = 0.05 and Holm-Bonferroni correction for family-wise error rate control.

\subsection{H1: Message Loss (CA2 confirmatory)}

\textbf{Source of truth:} \texttt{results/summary/stats\_H1\_H2\_H3.json} and \texttt{hypotheses.md} (\texttt{backend: localsim} only).

\textbf{Registered H1 (campaign A\_vt\_consumer\_kill):} visibility timeout has no effect on \emph{message-loss rate} under consumer failure.

\textbf{Result:} Loss rate was identically 0.0 for all VT levels tested (30--600\,s; $n=10$ each). With zero variance the confirmatory test is undefined (\texttt{test: none (no variance)}; $p=1.0$, Holm-adjusted $p=1.0$).

\textbf{Decision:} \textbf{Fail to reject H0} --- CA2 H1 as a VT$\rightarrow$loss effect is \textbf{not estimable} on this simulator campaign because the queue+DLQ arm never loses messages under the tested horizon. Zero loss is an important descriptive finding, but it is \textbf{not} the same as a supported VT-effect hypothesis. Exploratory analyses instead show VT effects on \emph{recovery time} and \emph{duplicate rate}.

\subsection{H2: maxReceiveCount and recovery time}

\textbf{Registered H2 (campaign B\_mrc\_unhandled\_error):} retry limit (\texttt{maxReceiveCount}) has no effect on recovery time.

\textbf{Test:} Kruskal--Wallis on recovery\_time\_s across MRC $\in\{1,3,5,10\}$ ($n=10$ each).

\textbf{Result (committed):} $H=4.98$, raw $p=0.173$, Holm-adjusted $p=0.520$, $\varepsilon^2=0.128$. Group medians: 30\,s (MRC\,1), 35\,s (MRC\,3), 30\,s (MRC\,5), 35\,s (MRC\,10).

\textbf{Decision:} \textbf{Fail to reject H0} (do not claim ``H2 supported'' as a positive finding beyond non-significance). Exploratory DLQ capture vs MRC on the same campaign is strong ($H=37.69$, $p=3.3\times10^{-8}$, $\varepsilon^2=0.966$): MRC\,1 mean DLQ capture $\approx0.294$ vs $\approx0$ at MRC$\ge5$.

\subsection{H3: Steady-state optimum under fault}

\textbf{Registered tests (campaign E\_guidance\_transfer):} throughput-optimal config (VT\,600, batch\,50) vs alternatives on (i)~loss and (ii)~recovery.

\textbf{H3\_loss:} loss identically 0 --- no variance; fail to reject H0 ($p_{\mathrm{adj}}=1$).

\textbf{H3\_recovery:} Mann--Whitney $U=250$, raw $p=7.92\times10^{-4}$, Holm-adjusted $p=3.17\times10^{-3}$, rank-biserial $r=0.667$. Optimal mean recovery $600$\,s vs alternatives mean $224$\,s ($n_{\mathrm{opt}}=10$, $n_{\mathrm{rest}}=30$).

\textbf{Decision:} \textbf{Reject H0 for H3\_recovery} (significant after Holm). Earlier draft text citing $p_{\mathrm{adj}}=0.084$ and ``borderline / not significant'' is \textbf{incorrect} relative to committed stats and is withdrawn. Practical reading: the steady-state-favourable long VT materially lengthens recovery under consumer kill in the simulator.
\section{Key Findings}

\subsection{Recovery Time Follows Visibility Timeout 1:1}

Exploratory campaign~A (\texttt{consumer\_kill}) recovery\_time\_s group means in \texttt{stats\_H1\_H2\_H3.json}:
\begin{itemize}
    \item VT\,30: mean $33$\,s (median $30$)
    \item VT\,60: mean $61$\,s
    \item VT\,120: mean $120$\,s
    \item VT\,300: mean $300$\,s
    \item VT\,600: mean $600$\,s
\end{itemize}
($H=48.50$, $p=7.44\times10^{-10}$, $\varepsilon^2=0.990$). Mechanism: failed messages remain in-flight until VT expires, then reappear for reprocessing.

\textbf{Implication:} \textbf{VT is a recovery time knob.} Set it based on acceptable recovery SLA, not just steady-state processing time.

\subsection{DLQ Traffic Depends on maxReceiveCount (exploratory)}

Committed exploratory Kruskal--Wallis results in \texttt{stats\_H1\_H2\_H3.json}:

\textbf{Campaign B (\texttt{unhandled\_error}), DLQ capture by MRC:}
\begin{itemize}
    \item MRC\,1: mean $\approx0.294$ ($n=10$)
    \item MRC\,3: mean $\approx0.0009$
    \item MRC\,5 and MRC\,10: mean $0.0$
    \item $H=37.69$, $p=3.28\times10^{-8}$, $\varepsilon^2=0.966$
\end{itemize}

\textbf{Campaign C (\texttt{datastore\_reject}), DLQ capture by MRC:}
\begin{itemize}
    \item MRC\,1: mean $\approx0.170$
    \item MRC\,3: mean $\approx0.0003$
    \item MRC\,5 and MRC\,10: mean $0.0$
    \item $H=32.70$, $p=3.72\times10^{-7}$, $\varepsilon^2=0.839$
\end{itemize}

\textbf{Implication:} Under these simulator fault windows, MRC\,1 dead-letters a large fraction that higher MRC retains for retry; MRC$\ge5$ sends almost nothing to DLQ in the committed cells. Do not cite older draft ranges (e.g., ``5--8\% DLQ at MRC$\ge5$'') that contradict these means.

\subsection{Duplicates Over Baseline Floor}

Baseline duplicate floor from H0 summary: $\approx0.00197$. Under \texttt{unhandled\_error} (campaign~B excess duplicates in the same JSON): MRC\,1 excess $\approx0$ (messages DLQ'd early); MRC\,3/5/10 excess $\approx0.136$ each. Higher MRC therefore trades DLQ traffic for duplicate redeliveries in this campaign. Idempotent consumers neutralize duplicate side effects; non-idempotent consumers do not.

\subsection{Batch Size and Recovery}

No confirmatory batch$\rightarrow$recovery test is registered in \texttt{stats\_H1\_H2\_H3.json}. Recovery under consumer kill is dominated by VT in exploratory campaign~A (means $\approx33$\,s at VT\,30 through $600$\,s at VT\,600). Batch-size recovery claims with $p=0.31$ / $\eta^2=0.03$ are withdrawn as unsupported by the committed stats files.

\section{Figures}

Eleven figures are provided in \texttt{results/figures/}:

\begin{enumerate}
    \item \texttt{loss\_vs\_visibility.png}: Loss rate by VT and fault type (zero across all)
    \item \texttt{recovery\_vs\_visibility.png}: Recovery time vs. VT (linear, slope ~1)
    \item \texttt{recovery\_vs\_max\_receive\_count.png}: Recovery time vs. MRC (flat for transient faults)
    \item \texttt{dlq\_heatmap.png}: DLQ rate heatmap (MRC × fault type)
    \item \texttt{dlq\_vs\_max\_receive\_count.png}: DLQ rate by MRC (bar chart)
    \item \texttt{duplicates\_over\_floor.png}: Duplicate rate increase by MRC
    \item \texttt{timeline\_visible\_vs\_backlog.png}: Queue depth timeline during fault (visible vs. total)
    \item \texttt{arms\_by\_fault.png}: Sync vs. queue arm metrics by fault type
    \item \texttt{baseline\_throughput\_latency.png}: Steady-state throughput/latency by batch size
    \item \texttt{guidance\_transfer.png}: Throughput vs. recovery time scatter (H3 configs)
    \item \texttt{baseline/baseline\_throughput\_latency.png}: Detailed baseline replication
\end{enumerate}

All figures include axis labels, legends, titles, and captions in report text.

\section{Summary}

The evaluation confirms:
\begin{itemize}
    \item \textbf{Queue-arm loss floor is zero} under tested faults (H1 VT$\rightarrow$loss confirmatory test not estimable).
    \item \textbf{MRC$\rightarrow$recovery H2:} fail to reject; DLQ capture vs MRC is the clear exploratory signal.
    \item \textbf{Long VT delays recovery} (H3\_recovery reject H0 after Holm, $p_{\mathrm{adj}}=0.003$, $r=0.667$).
    \item \textbf{Visibility timeout is the primary recovery time determinant} (1:1 relationship).
    \item \textbf{MRC should be tuned using DLQ/duplicate evidence} (exploratory: MRC\,1 high DLQ; MRC$\ge5$ near-zero DLQ but higher duplicates under unhandled\_error).
    \item \textbf{Batch$\rightarrow$recovery:} not supported by committed confirmatory stats; optimize batch for throughput with caution.
\end{itemize}

Findings are \textbf{localsim-only}. Adaptive VT/DIVE was not evaluated (\texttt{adaptive\_vt: false}). Cost was not an experimental DV. Live AWS validation remains required for CA2 cloud evidence.
